% Standalone problem document, generated from the adjacent README.md.
% Compile with XeLaTeX. Mathematical content is maintained in Markdown.
\documentclass[11pt,a4paper]{article}
\usepackage[fontsize=10.5pt]{scrextend}
\usepackage[margin=22mm,headheight=15pt,headsep=9mm,footskip=11mm]{geometry}
\usepackage{fontspec}
\setmainfont{texgyrepagella}[Extension=.otf,UprightFont=*-regular,BoldFont=*-bold,ItalicFont=*-italic,BoldItalicFont=*-bolditalic]
\setsansfont{texgyreheros}[Extension=.otf,UprightFont=*-regular,BoldFont=*-bold,ItalicFont=*-italic,BoldItalicFont=*-bolditalic]
\setmonofont{texgyrecursor}[Extension=.otf,UprightFont=*-regular,BoldFont=*-bold,ItalicFont=*-italic,BoldItalicFont=*-bolditalic,Scale=MatchLowercase]
\usepackage{amsmath,amssymb}
\usepackage{unicode-math}
\setmathfont{texgyrepagella-math.otf}
\usepackage{xcolor,microtype,enumitem,needspace,fancyhdr}
\usepackage{bookmark}
\definecolor{ink}{HTML}{17344B}
\definecolor{accent}{HTML}{197D80}
\definecolor{muted}{HTML}{596773}
\hypersetup{colorlinks=true,linkcolor=accent,urlcolor=accent,pdftitle={RE-01 - Constant-factor
HSS approximation in polynomial
time},pdfauthor={Open Problems in Numerical Linear Algebra}}
\urlstyle{same}
\setlength{\parindent}{0pt}
\setlength{\parskip}{5pt plus 1pt minus 1pt}
\linespread{1.04}
\setlength{\emergencystretch}{3em}
\widowpenalty=10000
\clubpenalty=10000
\displaywidowpenalty=10000
\allowdisplaybreaks[1]
\setlist{nosep,leftmargin=1.5em}
\providecommand{\tightlist}{\setlength{\itemsep}{0pt}\setlength{\parskip}{0pt}}
\providecommand{\pandocbounded}[1]{#1}
\pagestyle{fancy}
\fancyhf{}
\fancyhead[L]{\sffamily\footnotesize\color{muted}OPEN PROBLEMS IN NUMERICAL LINEAR ALGEBRA}
\fancyhead[R]{\sffamily\bfseries\footnotesize\color{accent}RE-01}
\fancyfoot[L]{\parbox[b]{0.9\textwidth}{\sffamily\fontsize{8}{10}\selectfont\color{muted}Editorial ratings; a bounded search cannot certify that a problem remains unsolved.\\Literature check: 2026-09-10}}
\fancyfoot[R]{\sffamily\footnotesize\color{muted}\thepage}
\renewcommand{\headrulewidth}{0.3pt}
\renewcommand{\headrule}{\hbox to\headwidth{\color{accent}\leaders\hrule height \headrulewidth\hfill}}
\makeatletter
\renewcommand{\section}{\Needspace{5\baselineskip}\@startsection{section}{1}{0pt}{12pt plus 2pt minus 2pt}{4pt}{\sffamily\bfseries\large\color{ink}}}
\renewcommand{\subsection}{\Needspace{4\baselineskip}\@startsection{subsection}{2}{0pt}{9pt plus 1pt minus 1pt}{3pt}{\sffamily\bfseries\normalsize\color{ink}}}
\makeatother
\setcounter{secnumdepth}{0}
\begin{document}
{\sffamily\small\color{muted}Randomized and low-rank approximation\par}
\vspace{3pt}
{\raggedright\hyphenpenalty=10000\exhyphenpenalty=10000\sffamily\bfseries\fontsize{21}{24}\selectfont\color{ink}Constant-factor
HSS approximation in polynomial time\par}
\vspace{7pt}
{\sffamily\small\textbf{Difficulty:} challenging\quad\textbf{Importance:} interesting
to the community\par}
{\sffamily\small\bfseries\color{accent}Status: Open\par}
\vspace{4pt}
{\color{accent}\hrule height 0.8pt}
\vspace{6pt}
\textbf{Rating rationale:} Challenging because shared HSS bases couple
approximation choices across levels; community impact is reliable
compression for hierarchical numerical solvers.

\textbf{Area:} structured low-rank matrix approximation

\subsection{Context and notation}\label{context-and-notation}

Use exact real arithmetic, with comparisons, standard Gaussian sampling,
and exact SVDs available as primitives; charge an SVD of an
\(a\times b\) matrix \(O(ab\min(a,b))\) operations. This states the
idealized arithmetic model used here, rather than a finite precision or
bit complexity claim. A matrix--vector query returns either \(Av\) or
\(A^\mathsf Tv\) for one chosen real vector \(v\); both types count
toward the total. Randomized guarantees are for every fixed input, with
probability or expectation over the algorithm's randomness.

For \(N=2^{L+1}k\), \(k,L\ge1\), form the perfect binary tree that
successively halves the ordered index set \([N]=\{1,\ldots,N\}\) until
every leaf contains \(2k\) indices. Define \(\mathcal S_{L,k}\) to
contain precisely the real \(N\times N\) matrices \(B\) such that, for
every nonroot node \(I\) of this tree,

\[
\operatorname{rank}B[I,[N]\setminus I]\le k,
\qquad
\operatorname{rank}B[[N]\setminus I,I]\le k.
\]

This is the HSS class with rank at most \(k\), including its constraints
across levels. There is no restriction inside a leaf diagonal block. Set

\[
E_{L,k}(A)=\min_{B\in\mathcal S_{L,k}}\|A-B\|_F^2.
\]

The minimum exists. The rank definition follows Stefano Massei, Leonardo
Robol, and Daniel Kressner,
\href{https://arxiv.org/pdf/1909.07909v3}{\emph{hm-toolbox: MATLAB
Software for HODLR and HSS Matrices}}, SIAM Journal on Scientific
Computing \textbf{42} (2020), C43--C68, §2.2, Definition 3. Equivalence
with the telescoping representation and existence of a minimizer are
recorded in Amsel et al., cited below, Definition 3 and Appendix D.

\subsection{Problem statement}\label{problem-statement}

Do absolute constants \(C,c>0\) and a uniform randomized algorithm exist
that, given the entries of any \(A\in\mathbb R^{N\times N}\) and
integers \(k,L\ge1\) with \(N=2^{L+1}k\), returns
\(B\in\mathcal S_{L,k}\) using \(O(N^c)\) arithmetic operations and
satisfies

\[
\mathbb E\|A-B\|_F^2\le C\,E_{L,k}(A)?
\]

The exponent and approximation constant must be independent of
\(k,L,A\). A deterministic algorithm is allowed. The output must retain
rank at most \(k\).

\newpage
\subsection{Reference}\label{reference}

Noah Amsel, Tyler Chen, Feyza Duman Keles, Diana Halikias, Cameron
Musco, Christopher Musco, and David Persson,
\href{https://arxiv.org/html/2505.16937v2}{\emph{Quasi-optimal
Hierarchically Semi-separable Matrix Approximation}}, SIAM Journal on
Matrix Analysis and Applications \textbf{47} (2026), 586--621, §3.4;
Theorems 6 and 12. That section expressly asks for a polynomial-time
constant-factor approximation.

\subsection{Status evidence}\label{status-evidence}

The known squared-error factor is \(O(L)\). The lower bound near \(2\)
concerns the analyzed greedy algorithm, not all algorithms. Searches for
``HSS constant-factor approximation 2026'' and ``hierarchically
semiseparable approximation constant 2026'' found no resolution. The
latest arXiv version is v2, September 6, 2025; the
\href{https://doi.org/10.1137/25M176622X}{journal version} appeared
online May 8, 2026.

\subsection{Audit --- 2026-09-10}\label{audit-2026-09-10}

Rechecked \href{https://arxiv.org/html/2505.16937v2}{Amsel et al., §3.4
and Appendix B.1} and its \href{https://doi.org/10.1137/25M176622X}{2026
journal record}. Constant-factor polynomial-time HSS approximation
remains unresolved. Constant-factor and later HSS searches found no
solution; the known depth-dependent factor is weaker than this target.
\end{document}
